\chapter{Discussion and Conclusion}
\label{sec:conclusion}

% **************************** Define Graphics Path **************************
\ifpdf
    \graphicspath{{Chapter10/Figs/Raster/}{Chapter10/Figs/PDF/}{Chapter10/Figs/}}
\else
    \graphicspath{{Chapter10/Figs/Vector/}{Chapter10/Figs/}}
\fi

This final chapter presents a comprehensive summary of the findings detailed in this thesis, alongside an acknowledgment of its limitations and proposals for future research directions.

\section{Discussion}
\label{ch3:sec:disc}
In the first part of this thesis, we explored MTL and DA techniques aimed at improving the quality of ML in binary image classification tasks under underrepresentation constraints. The methodologies presented—ranging from MTL and OS to ADV and diffusion-model-based approaches—proved to be effective techniques for mitigating cultural bias. These methods can be integrated into a comprehensive pipeline for both pre-processing and in-processing mitigation strategies, as demonstrated by the MTL+PARB+DM case.

However, some of these techniques present limitations regarding their application or computational requirements. RT and diffusion-model-based techniques are highly task-dependent; while some generalization to other domains is possible, its implementation is non-trivial.

Furthermore, diffusion-model-based and ADV techniques demand a significant amount of time and resources, which may not be practically justified for a small CIC, where MTL or OS techniques alone might suffice.

The MTL strategy can also be memory-intensive in cases involving multiple cultures or large "heads." In the latter scenario, a potential solution involves phased training: starting with one head (e.g., the majority culture) and subsequently training others. This approach limits peak memory consumption by shifting the bottleneck toward computational complexity.

The final section of this thesis introduces a LLM-based cognitive architecture engineered to embed cultural competence within robotic control loops. This architectural framework can be deployed under two operational paradigms: a static configuration, which instantiates a generalized or nation-specific cultural framework, or a dynamic configuration, which initializes a national cultural default as an a priori baseline and subsequently adapts via real-time user feedback. To rigorously evaluate the platform, cultural competence was decomposed into distinct metrics for "culture" and "competence," measuring each parameter independently before synthesizing them into an aggregate evaluation framework.

Empirical results demonstrate that while static architectures are highly rated by participants under specific conditions, they exhibit several operational constraints. Specifically, a culturally neutral baseline architecture is generally perceived as more competent overall, whereas a culture-specific architecture is evaluated as being more relatable and socially proximate to human interactants. Furthermore, executing real-time user adaptation remains a significant challenge for LLMs; the models frequently display latency and consistency bottlenecks when dynamically shifting linguistic and behavioral patterns at runtime. These limitations are further pronounced in quantized variants, which exhibit diminished performance compared to their full-precision counterparts.

\section{Limitations and Future Directions}

The experiments conducted in the first phase of this research are bounded in both domain and task scope. While the current implementation focuses on specific vision tasks, ML models can address a broad spectrum of culture-specific problems within the image domain, such as object detection and semantic segmentation. For these broader visual applications, the deployment of multi-task learning (MTL) and data augmentation (DA) algorithms remains highly valid.

When extending this paradigm to alternative modalities, natural language stands out as the primary vehicle for cultural transmission. This transmission occurs both explicitly, through vocabulary, dialect, and communication styles, and implicitly, through embedded behavioral norms. For instance, a system recommending a dinner reservation at 8:00 PM may be perceived as unusually late in certain geographic regions, yet exceptionally early in others, illustrating the latent cultural variables inherent in linguistic interaction.

Further expansions of MTL architecture and CIC metrics could be explored, relaxing the hard constraint of national categorization and using approximations of distribution-based metrics (which are more complex to compute, but may gather more fine-graded underrepresented cultures inside the main ones).

Regarding the final investigation of this dissertation, because these experiments were conceived as exploratory probes into the intersection of cultural competence and emergent ML technologies, several promising future research vectors emerge. First, conducting a systematic comparative analysis across a broader spectrum of LLMs—specifically evaluating their adherence to dense prompt constraints—would significantly reinforce the empirical validity of the core architectural assertions. Furthermore, systematically varying contextual and environmental variables (such as spatial layouts, robotic hardware configurations, targeted cultural demographics, localized languages, and conversational domains) would extend the generalizability of these foundational insights.

Critically, this research highlighted significant performance disparities in current ML models when evaluated across different national cultures. In alignment with the data constraints discussed in the vision-centric phase of this work, this phenomenon is primarily attributable to severe training data imbalances. The relative underrepresentation of Italian and German corpora within foundational training datasets, compared to the ubiquity of English-language data, directly accounts for why English-language interactions were generally perceived by participants as more fluid and structurally competent.

\section{Conclusion}

This dissertation systematically investigated the critical impact of cultural underrepresentation within ML models deployed in social robotics, developing and evaluating multi-tiered strategies to mitigate these systemic biases. The first phase of this research focused on algorithmic interventions within robotic vision pipelines, demonstrating that cultural biases must be addressed via specialized, task-specific ML frameworks. Without these ad hoc solutions, systematic representational errors inevitably perpetuate, resulting in ethically compromised and operationalized failures during human-robot deployments.

The second phase of this work served a dual purpose. First, it introduced a novel, exploratory LLM-based cognitive architecture engineered to navigate national cultural frameworks through common-sense reasoning and runtime behavioral adaptation. Second, it exposed the structural boundaries of contemporary generative algorithms, which remain severely constrained in their ability to represent nuanced national and sub-national cultures without reverting to superficial stereotypes. When interacting with end-users, these models frequently exhibit operational volatility due to the foundational data imbalances identified in the early stages of this work, reinforcing the absolute necessity of robust structural mitigation strategies.

Crucially, the empirical evaluation of this adaptation strategy revealed a critical systemic threshold in human-robot cross-cultural interaction. While the online calibration loop and LLM-driven adaptation were theoretically engineered to mitigate sub-national data scarcity, their perceived utility was strictly bounded by fundamental operational competence. When low-level task execution or natural language understanding faltered, high-level cultural adaptation strategies failed to yield statistically significant improvements in user engagement, and in several cohorts, exacerbated user frustration. This outcome establishes a crucial design axiom for culturally competent social robotics: high-level behavioral alignment cannot compensate for underlying operational volatility. Consequently, low-level programmatic constraint enforcement—such as the framework developed in the first part of this thesis—must be treated as an absolute prerequisite before online behavioral adaptation algorithms can be effectively deployed.

This research contributes a formal framework for operationalizing ``cultural competence'' within the intersection of ML and social robotics. Nevertheless, a vast landscape remains for future academic development. The phenomenon of Cultural Incompetence requires rigorous future study regarding its distinct manifestations across alternative modalities and environmental domains. Furthermore, diverse methodologies and larger-scale human-subject studies are necessary to complement these findings and expand the existing literature. It remains evident that the concept of ``culture'' still lacks a cross-disciplinary, formal operational definition, which continues to fragment both the evaluation methodologies and data analyses within the broader field.

Consequently, this thesis argues that distribution-based approaches—which formalize culture through fine-grained, non-categorical statistical densities—could significantly advance a research domain that has historically remained operationally stagnant by conceptualizing culture solely as a rigid national construct. While the proposed CIC metric can be dynamically adjusted to accommodate these granular cultural definitions, transitioning toward entirely non-categorical, fluid representations remains a significant open challenge for the near future. 

Crucially, this research highlights the pronounced degree of behavioral and linguistic bias remaining in edge-deployed, quantized LLMs. Despite their accelerating integration into commercial robotic platforms, these architectures continue to exhibit severe bias and performance degradation even when evaluating well-documented national cultures, demonstrating that post-training quantization disproportionately strips away subtle cultural alignment vectors.

Ultimately, future research must prioritize empirical, human-in-the-loop experimentation to prevent the real-world deployment of unethical or malfunctioning SRs. While discussions surrounding AI ethics and cultural alignment have historically remained largely theoretical, the human-subject research presented in this dissertation demonstrates that participants provide vital diagnostic insights into systematic errors that are otherwise undetectable in simulation. As a bridge between behavioral sciences and robotic engineering, culturally competent social robotics must shift its paradigm toward investigating not only idealized system successes, but specifically the mechanics of what is not working and why.